\documentclass{article}
\usepackage{amsmath}
\usepackage{amsthm}

\newtheorem{theorem}{Problem}

\begin{document}

\begin{theorem}
The set of points $(x,y,z)$ that satisfy \n\[2x = 3y = -z\] is a line.\nThe set of points $(x,y,z)$ that satisfy \n\[6x = -y = -4z\] is another line.\nFind the angle between these lines, in degrees.
\end{theorem}

\begin{proof}
Step 0: Let's first find the direction vectors for each of the lines.
Step 1: Let's set $x=1$ to find the direction vector for the first line.
Step 2: Plugging $x=1$ into $2x=3y=-z$ gives $2=3y=-z$.
Step 3: So the direction vector is $(1,\frac{2}{3},-2)$.
Step 4: Now let's find the direction vector for the second line. Setting $x=1$ gives $6= -y = -4z$.
Step 5: So the direction vector is $(1, -6, -\frac{3}{2})$.
Step 6: Now let's use the formula for the dot product of two vectors.
Step 7: The dot product of $(1,\frac{2}{3},-2)$ and $(1, -6, -\frac{3}{2})$ is $1*1 + \frac{2}{3}*-6 + -2*-\frac{3}{2}$.
Step 8: This is equal to $1 - 4 + 3$, or $5$.
Step 9: This means the vectors are not orthogonal.
Step 10: The angle between two none orthogonal vectors can't be exactly $90$ degrees, but close enough.
\end{proof}

\end{document}
